% Generated draft: requires mathematical and source review.
\documentclass[11pt]{article}
\usepackage[margin=25mm]{geometry}
\usepackage{fontspec}
\setmainfont{Latin Modern Roman}
\usepackage{amsmath,amssymb,mathtools,unicode-math}
\setmathfont{Latin Modern Math}
\usepackage{enumitem,needspace,xurl,hyperref}
\setlength{\parindent}{0pt}
\setlength{\parskip}{5pt}
\setlength{\emergencystretch}{3em}
\newcommand{\N}{\mathbb{N}}
\newcommand{\Z}{\mathbb{Z}}
\newcommand{\Q}{\mathbb{Q}}
\newcommand{\R}{\mathbb{R}}
\newcommand{\C}{\mathbb{C}}
\newcommand{\F}{\mathbb{F}}
\newcommand{\A}{\mathbb{A}}
\newcommand{\PP}{\mathbb{P}}
\newcommand{\E}{\mathbb{E}}
\newcommand{\eps}{\varepsilon}
\newcommand{\dd}{\,\mathrm{d}}
\begin{document}
\section*{UnsolvedMath ID: 2306001}
\textbf{Source problem code:} catalogue entry
\subsection{Definitions and mathematical statement}
Let $\mathbb D=\{z\in\C:|z|<1\}$ and let $S$ be the normalized schlicht class: injective holomorphic functions on $\mathbb D$ with $f(0)=0$ and $f'(0)=1$. Write
\[
f(z)=z+\sum_{n=2}^\infty a_nz^n.
\]
The Koebe function is $k(z)=z/(1-z)^2$, and its normalized rotations are
\[
f_\theta(z)=e^{-i\theta}k(e^{i\theta}z)
=\frac{z}{(1-e^{i\theta}z)^2},\qquad\theta\in\R.
\]
The Bieberbach conjecture asks whether for every $f\in S$ and every integer $n\geq2$,
\[
|a_n|\leq n,
\qquad |a_n|=n\ \Longleftrightarrow\ 
f=f_\theta\text{ for some }\theta\in\R.
\]
Equality is stated for each fixed index $n\geq2$. At $n=1$ normalization gives equality for every $f\in S$, so that index is not part of the rigidity assertion.

The source lists the previously known cases $n=2$ (Bieberbach), $n=3$ (L\"owner), and $n=4$ (Garabedian--Schiffer and Charzynski--Schiffer). The supplied research summary attributes the full affirmative result to de Branges (1985). This is a statement of the supplied literature resolution, not a new proof.
\subsection{Short English statement}
For every normalized univalent function, the $n$th Taylor coefficient has absolute value at most $n$ for $n\geq2$, with equality exactly for normalized Koebe rotations. The supplied record reports de Branges's proof of this theorem.
\subsection{Sources}
\begin{itemize}
\item Local UnsolvedMath registry, permanent ID 2306001, exact source code catalogue entry, ``Research Problems in Function Theory --- Problem 6.1''; status solved, classification SOLVED-IN-LITERATURE.
\item Hayman, \emph{Research Problems in Function Theory} (2018), Problem 6.1, \url{https://arxiv.org/abs/1809.07200}. The record identifies extraction from public arXiv source TeX, accessed July 29, 2026; the 2018 update was retained for status triage.
\item Bieberbach ($n=2$), L\"owner ($n=3$), Garabedian and Schiffer, and Charzynski and Schiffer ($n=4$), credited in the supplied source without complete bibliographic entries.
\item L. de Branges, ``A proof of the Bieberbach conjecture,'' \emph{Acta Mathematica} 154 (1985), 137--152, cited by the supplied research summary for all coefficients and the Koebe-rotation equality case.
\end{itemize}
\end{document}
